% Unit 1: Foundations and Prerequisites
% Practice Problems and Solutions (Part II and Part III only)
% Prepared by Staples Education
% Content informed by the local curriculum modules (1.1 to 1.6) and the
% QUIZ_DATA micro-practice and unit-mastery items for Unit 1.

\documentclass[11pt]{article}

\usepackage[margin=1in]{geometry}
\usepackage{amsmath}
\usepackage{amssymb}
\usepackage{xcolor}
\usepackage[most]{tcolorbox}
\usepackage{enumitem}
\usepackage{hyperref}
\hypersetup{colorlinks=true, linkcolor=headerblue, urlcolor=headerblue}

\tcbuselibrary{skins}

% ---------------------------------------------------------------------------
% Unified style-guide palette (see LATEX_STYLE_GUIDE.md)
% ---------------------------------------------------------------------------
\definecolor{headerblue}{HTML}{1C569E}   % overview / structural framing
\definecolor{accentgreen}{HTML}{2E7D32}  % algorithms / methods
\definecolor{accentorange}{HTML}{E27A1E} % formulas / concept takeaways
\definecolor{workpurple}{HTML}{A020F0}   % worked-example tracking (vibrant)
\definecolor{warnred}{HTML}{C0392B}      % crimson pitfall / warning

% ---------------------------------------------------------------------------
% Subsection typography: headerblue numbered sub-module titles
% ---------------------------------------------------------------------------
\usepackage{titlesec}
\titleformat{\subsection}
  {\normalfont\large\bfseries\color{headerblue}}{\thesubsection}{1em}{}

% ---------------------------------------------------------------------------
% Six core custom boxes (unbreakable, title={#1} protected)
% ---------------------------------------------------------------------------

% 1. Blue overview capsule for the topics summary.
\newtcolorbox{topicsbox}{
    enhanced,
    colback=headerblue!6!white, colframe=headerblue,
    coltitle=white, fonttitle=\bfseries,
    title=Unit 1 at a Glance: Core Sub-Topics,
    boxrule=0.6pt, arc=2mm,
    left=3mm, right=3mm, top=2mm, bottom=2mm
}

% 2. Orange formula container for standard forms and master formulas.
\newtcolorbox{formulabox}[1][Master Formula]{
    enhanced,
    colback=accentorange!8!white, colframe=accentorange,
    coltitle=white, fonttitle=\bfseries,
    title={#1},
    boxrule=0.6pt, arc=1mm,
    left=3mm, right=3mm, top=2mm, bottom=2mm
}

% 3. Green algorithm block for explicitly laid-out procedures.
\newtcolorbox{methodbox}[1][Method]{
    enhanced,
    colback=accentgreen!6!white, colframe=accentgreen,
    coltitle=white, fonttitle=\bfseries,
    title={#1},
    boxrule=0.6pt, arc=1mm,
    left=3mm, right=3mm, top=2mm, bottom=2mm
}

% 4. Purple west-bordered worked example for step-by-step solutions.
\newtcolorbox{examplebox}[1][Worked Example]{
    enhanced,
    colback=workpurple!4!white, colframe=workpurple,
    coltitle=white, fonttitle=\bfseries,
    title={#1},
    boxrule=0.4pt, sharp corners,
    borderline west={3pt}{0pt}{workpurple},
    left=5mm, right=3mm, top=2mm, bottom=2mm
}

% 5. Orange thin-bordered final concept takeaway.
\newtcolorbox{conceptbox}[1][Key Takeaway]{
    enhanced,
    colback=accentorange!5!white, colframe=accentorange,
    coltitle=white, fonttitle=\bfseries,
    title={#1},
    boxrule=0.3pt, arc=1mm,
    left=3mm, right=3mm, top=2mm, bottom=2mm
}

% 6. Crimson west-bordered warning box for common pitfalls.
\newtcolorbox{warningbox}[1][Common Pitfall]{
    enhanced,
    colback=red!3!white, colframe=warnred,
    coltitle=white, fonttitle=\bfseries,
    title={#1},
    boxrule=0.4pt, sharp corners,
    borderline west={3pt}{0pt}{warnred},
    left=5mm, right=3mm, top=2mm, bottom=2mm
}

\setlist[enumerate]{itemsep=4pt, topsep=4pt}

% Number and index sections plus subsections (e.g., 1.1) two levels deep.
\setcounter{secnumdepth}{2}
\setcounter{tocdepth}{2}

\begin{document}

% ---------------------------------------------------------------------------
% Title block
% ---------------------------------------------------------------------------
\begin{center}
    {\LARGE \bfseries Unit 1: Foundations and Prerequisites}\\[6pt]
    {\large Prepared by Staples Education}\\[4pt]
    {\itshape Practice Problems and Solutions}
\end{center}

\vspace{0.6em}
\hrule
\vspace{0.6em}

% ===========================================================================
% PART II: REVIEW GUIDE (PRACTICE QUESTIONS)
% No \clearpage here: a standalone Practice Set opens directly on its problems.
% ===========================================================================
\section*{Part II --- Review Guide: Practice Problems}

Work the following ten problems in order. They are graded from raw calculus
review, through the number $e$ and Euler's formula, into classifying equations,
verifying solutions, and finally building equations by eliminating constants.
Complete solutions follow in Part III.

\begin{enumerate}[leftmargin=2em]
    \item Differentiate $y = e^{3x}$ using the chain rule.

    \item Evaluate the indefinite integral $\displaystyle\int e^{-2x}\,dx$.

    \item Differentiate \textbf{(a)} $(2x+1)^{3}$ and \textbf{(b)} $\sin(x^{2})$.

    \item Compute $\dfrac{d}{dt}e^{kt}$, then state the general solution of
          $\dfrac{dy}{dt} = y$.

    \item Use Euler's formula $e^{i\theta} = \cos\theta + i\sin\theta$ to
          evaluate \textbf{(a)} $e^{i\pi}$ and \textbf{(b)} $e^{i\pi/2}$, and
          write down Euler's identity.

    \item For each equation, state the order and the degree, and say whether it
          is linear:
          \textbf{(a)} $\dfrac{d^{2}y}{dx^{2}} + 3\dfrac{dy}{dx} = 0$, \quad
          \textbf{(b)} $\left(\dfrac{dy}{dx}\right)^{3} + y = x$, \quad
          \textbf{(c)} $\dfrac{dy}{dx} + y^{2} = 0$.

    \item Verify that $y = e^{2x}$ is a solution of $\dfrac{d^{2}y}{dx^{2}} - 4y = 0$.

    \item Decide whether $y = x^{2}$ satisfies $x\dfrac{dy}{dx} = 2y$, and
          whether $y = \cos x$ satisfies $\dfrac{d^{2}y}{dx^{2}} + y = 0$.

    \item Form a differential equation by eliminating $C$ from \textbf{(a)}
          $y = Cx^{2}$ and \textbf{(b)} $y = Ce^{x}$.

    \item \textbf{(Capstone)} Eliminate both constants from
          $y = A\cos x + B\sin x$, and separately eliminate $r$ from the circle
          family $x^{2} + y^{2} = r^{2}$.
\end{enumerate}

\vspace{0.4em}
\hrule
\vspace{0.6em}

% ===========================================================================
% PART III: SOLUTIONS
% \clearpage so the solutions print on their own page.
% ===========================================================================
\clearpage
\section*{Part III --- Complete Solutions}

\noindent\textbf{Solution 1.}\quad The exponential reproduces itself, and the
chain rule multiplies by the derivative of the inside, $3$.
\[ \textcolor{workpurple}{\frac{dy}{dx} = e^{3x}\cdot 3 = 3e^{3x}}. \]

\medskip
\noindent\textbf{Solution 2.}\quad Integrating an exponential divides by the
constant in the exponent and restores the arbitrary constant.
\[ \textcolor{workpurple}{\int e^{-2x}\,dx = \frac{1}{-2}e^{-2x} + C = -\tfrac{1}{2}e^{-2x} + C}. \]

\medskip
\noindent\textbf{Solution 3.}\quad Apply the power-and-chain rule, then the
chain rule on the sine.
\[ \textcolor{workpurple}{\frac{d}{dx}(2x+1)^{3} = 3(2x+1)^{2}\cdot 2 = 6(2x+1)^{2}} \]
\[ \textcolor{workpurple}{\frac{d}{dx}\sin(x^{2}) = \cos(x^{2})\cdot 2x = 2x\cos(x^{2})}. \]

\medskip
\noindent\textbf{Solution 4.}\quad The exponential returns itself times the
constant from the exponent.
\[ \textcolor{workpurple}{\frac{d}{dt}e^{kt} = k\,e^{kt}}. \]
The equation $\frac{dy}{dt} = y$ asks for a function equal to its own derivative,
so
\[ \textcolor{workpurple}{y = Ce^{t}}. \]

\medskip
\noindent\textbf{Solution 5.}\quad Substitute the angles into Euler's formula.
\[ \textcolor{workpurple}{e^{i\pi} = \cos\pi + i\sin\pi = -1 + 0 = -1} \]
\[ \textcolor{workpurple}{e^{i\pi/2} = \cos\tfrac{\pi}{2} + i\sin\tfrac{\pi}{2} = 0 + i = i}. \]
Rearranging the first gives Euler's identity,
$\textcolor{workpurple}{e^{i\pi} + 1 = 0}$.

\medskip
\noindent\textbf{Solution 6.}\quad The order is the highest derivative; the
degree is its power; linearity forbids powers above one on $y$ and its
derivatives.
\[ \textcolor{workpurple}{\text{(a) order } 2,\; \text{degree } 1,\; \text{linear}} \]
\[ \textcolor{workpurple}{\text{(b) order } 1,\; \text{degree } 3,\; \text{nonlinear}} \]
\[ \textcolor{workpurple}{\text{(c) order } 1,\; \text{degree } 1,\; \text{nonlinear (the } y^{2} \text{ term)}}. \]

\medskip
\noindent\textbf{Solution 7.}\quad Differentiate twice and substitute.
\[ \textcolor{workpurple}{\frac{dy}{dx} = 2e^{2x}}, \qquad
   \textcolor{workpurple}{\frac{d^{2}y}{dx^{2}} = 4e^{2x}} \]
\[ \textcolor{workpurple}{\frac{d^{2}y}{dx^{2}} - 4y = 4e^{2x} - 4e^{2x} = 0} \qquad\checkmark \]
so $y = e^{2x}$ is a solution.

\medskip
\noindent\textbf{Solution 8.}\quad Test each candidate by substitution.
\[ \textcolor{workpurple}{y = x^{2}: \quad y' = 2x, \;\; x\,y' = x\cdot 2x = 2x^{2} = 2y} \qquad\checkmark \]
\[ \textcolor{workpurple}{y = \cos x: \quad y'' = -\cos x, \;\; y'' + y = -\cos x + \cos x = 0} \qquad\checkmark \]
Both candidates check out.

\medskip
\noindent\textbf{Solution 9.}\quad Differentiate, then replace the constant from
the original relation.
\[ \textcolor{workpurple}{y = Cx^{2}: \quad y' = 2Cx, \;\; C = \frac{y}{x^{2}} \;\Longrightarrow\; y' = 2\cdot\frac{y}{x^{2}}\cdot x = \frac{2y}{x}} \]
\[ \textcolor{workpurple}{\text{so } x\frac{dy}{dx} = 2y}. \]
\[ \textcolor{workpurple}{y = Ce^{x}: \quad y' = Ce^{x} = y \;\Longrightarrow\; \frac{dy}{dx} = y}. \]

\medskip
\noindent\textbf{Solution 10.}\quad Each family needs differentiation to clear
its constants. For $y = A\cos x + B\sin x$ (two constants), differentiate twice:
\[ \textcolor{workpurple}{y'' = -A\cos x - B\sin x = -y \;\Longrightarrow\; \frac{d^{2}y}{dx^{2}} + y = 0}. \]
For the circle $x^{2} + y^{2} = r^{2}$ (one constant), differentiate implicitly:
\[ \textcolor{workpurple}{2x + 2y\frac{dy}{dx} = 0 \;\Longrightarrow\; \frac{dy}{dx} = -\frac{x}{y}}. \]

\vspace{0.8em}
\begin{conceptbox}[Unit 1 Key Takeaway]
These problems exercise the two directions of the subject. Going forward, you
\textcolor{accentorange}{differentiate and integrate} fluently and
\textcolor{accentorange}{verify a solution} by substituting it into the equation;
going backward, you \textcolor{accentorange}{eliminate arbitrary constants} to
build the equation a family satisfies. Throughout, the
\textcolor{accentorange}{order, degree, and linearity} of an equation are the
first labels you read off.
\end{conceptbox}

\end{document}
